
\documentclass[11pt,a4paper]{article}
\usepackage[margin=20mm]{geometry}
\usepackage[T1]{fontenc}
\usepackage{lmodern,tabularx,booktabs,hyperref,xcolor,graphicx,float}
\usepackage{enumitem}
\setlength{\parindent}{0pt}
\setlength{\parskip}{4pt}
\renewcommand{\arraystretch}{1.08}
\hypersetup{colorlinks=true,linkcolor=blue,urlcolor=blue}
\setlist{nosep,leftmargin=*}
\newcommand{\code}[1]{\texttt{\detokenize{#1}}}

\begin{document}

\begin{center}
{\Large\textbf{23CSE455 Design Patterns}}\\[-1pt]
{\LARGE\textbf{Case Study Report}}\\[2pt]
\textbf{Design Pattern Identification and Evolution in a React and Node.js-Based Restaurant Management Application}
\end{center}

\textbf{Team:} B1 \hfill \textbf{Date:} 11 October 2026

\textbf{Application:} Restaurant POS System \quad
\textbf{Team repository:} \href{https://github.com/SrivarmaBattini/DP_Project}{DP\_Project}

\textbf{Original repository:}
\href{https://github.com/iono-such-things/restaurant-pos-system}{restaurant-pos-system}

\textbf{Team members}

\begin{tabularx}{\linewidth}{@{}lllX@{}}
\toprule
Roll number & Student name & Batch & Contribution area\\
\midrule
AM.SC.U4AIE23032 & Surya Dingari & AIE-A & Facade; Chain of Responsibility\\
AM.SC.U4AIE23110 & Vijay Reddy & AIE-B & Strategy; Observer\\
AM.SC.U4AIE23057 & Srivarma & AIE-A & Adapter; State\\
AM.SC.U4AIE23039 & Bruhadesh & AIE-A & Factory Method; Decorator\\
AM.SC.U4AIE23053 & Teja & AIE-A & Command; Template Method\\
\bottomrule
\end{tabularx}

\textit{Members should confirm that the contribution areas reflect their completed work.}

\section{Application and architecture}

The selected application is a restaurant management and point-of-sale system. Its documented technology stack includes React and TypeScript on the frontend, Node.js and Express on the backend, Prisma for database access, and Socket.IO for real-time communication. The application supports restaurant workflows such as menu access, order creation, payment handling and kitchen updates.

\begin{figure}[H]
\centering
\includegraphics[width=0.90\linewidth,height=0.40\textheight,keepaspectratio]
{Restaurant POS Baseline Architecture.png}
\caption{High-level application architecture reconstructed from repository documentation. This is a conceptual representation, not a commit-verified source diagram.}
\end{figure}

\newpage
Requests travel from the React interface to Express routes, through middleware and application services. Services use Prisma for persistence and may publish real-time updates through Socket.IO.

\textbf{Baseline details:}
The inspected local Git baseline was commit accessed on 11 October 2026.
The original repository is available at
\url{https://github.com/iono-such-things/restaurant-pos-system}.
No root-level license file was found during inspection;
therefore, the repository's license status remains unverified.
The documented technology stack includes React, Express,
Prisma, Socket.IO and TypeScript.

\textbf{Verification note:} The audit recorded 28/28 local assertions passing. These assertions do not establish full-stack HTTP, database or live Socket.IO integration.

\section{Part 1: Pattern discovery}

Pattern candidates were identified by inspecting source responsibilities, dependencies, method calls and collaborations. Framework usage alone was not treated as proof of a GoF pattern.

\begin{tabularx}{\linewidth}{@{}lXlX@{}}
\toprule
Pattern & Candidate instance & Count & Status\\
\midrule
Facade & Authentication service & 1 & Confirmed\\
Facade & Order creation service & 1 & Confirmed\\
Observer-like & Restaurant/kitchen event collaboration & 1 & Confirmed as event-based collaboration\\
Chain of Responsibility & Express authentication middleware & 1 & Uncertain\\
Active Record & Prisma models & 1 & Rejected\\
\bottomrule
\end{tabularx}

\textbf{Authentication Facade:}
\code{AuthService.login} in
\code{apps/backend/src/services/auth.service.ts} coordinates authentication operations involving Prisma, bcrypt and JWT. It provides callers with one service entry point for a multi-step workflow.

\textbf{Order creation Facade:}
\code{OrderService.createOrder} in
\code{apps/backend/src/services/order.service.ts} coordinates order logic, persistence and event publication.

\textbf{Observer-like collaboration:}
\code{apps/backend/src/websocket.ts} and the frontend \code{KitchenDisplay.tsx} participate in real-time order and kitchen updates. This is described as Observer-like publish/subscribe behavior; the evidence does not establish a separate canonical GoF Subject/Observer class hierarchy.

\textbf{Uncertain/rejected candidates:}
Express middleware in \code{apps/backend/src/middleware/auth.ts} uses \code{next()}, but a custom handler hierarchy was not established. Therefore, Chain of Responsibility remains uncertain. Active Record is rejected because persistence is performed through Prisma Client rather than model-owned persistence methods.

The inspected baseline therefore supports two Facade instances and one Observer-like event collaboration. Source line ranges should be recorded for each classification, and the same classifications should be used consistently in the evidence files and final report.

\section{Part 2: Evolution and refactoring}

\textbf{CH01 -- Online ordering.}
Public menu and online-order routes allow customers to submit orders. The service validates input, calculates prices and publishes order events. The change extends the order-service workflow, but HTTP execution and database-backed order creation have not been verified end to end. A frontend preview fallback is not proof of a persisted order.

\textbf{CH02 -- Kitchen workflow.}
Orders are mapped to stations such as cold, hot and bar, with station-specific events published to connected clients. Nine local mapping assertions passed. Live Socket.IO delivery and station-room behavior remain unverified.

\textbf{CH03 -- Payment strategies.}
\code{IPaymentStrategy}, concrete strategies and \code{PaymentContext} separate payment selection from the caller. Local assertions cover strategy selection, but card and mobile implementations simulate success; real payment-provider integration is not established.

\begin{tabularx}{\linewidth}{@{}lXX@{}}
\toprule
Change & Design effect & Remaining verification\\
\midrule
CH01 & Extends order service and event flow & HTTP and database integration\\
CH02 & Adds station mapping and event routing & Live Socket.IO delivery\\
CH03 & Introduces Strategy-based payment selection & Persistence, failure handling and provider integration\\
\bottomrule
\end{tabularx}

The 28/28 local assertions are not full-stack integration tests. The evolved working tree was uncommitted, and complete behavior-preservation tests and before/after LOC or complexity measurements were not established. Record an immutable evolved commit and report only measured metrics. A refactoring should be described as behavior-preserving only when relevant tests support that claim.

\section{Part 3: Cross-language comparison}

Ten patterns were implemented as standalone demonstrations in Java, Python, JavaScript and C++: Adapter, Command, Chain of Responsibility, Decorator, Facade, Factory Method, Observer, State, Strategy and Template Method.

\begin{tabularx}{\linewidth}{@{}lrrX@{}}
\toprule
Language & Programs & Passed/verified & Result\\
\midrule
Java & 10 & 10 & Executed successfully\\
Python & 10 & 10 & Executed successfully\\
JavaScript & 10 & 10 & Executed successfully\\
C++ & 10 & 10 & Compiled and executed successfully\\
\bottomrule
\end{tabularx}

All 40 standalone programs were recorded as successfully executed: 10 programs in each of the four languages. The recorded environments included Node.js 22.21.1, CPython 3.10.11 and OpenJDK/javac 11.0.30.

Source LOC and timing evidence are in
\code{data/cross_language.csv} and
\code{Part3/evidence/reproduced_2026-10-11/}. The timing procedure used a warm-up and five process-time measurements per runnable example. These results characterize the supplied demonstrations and environment, not the languages in general. Maximum function cyclomatic complexity remains unavailable and should be measured with a consistent tool and settings before numerical comparisons are reported.

The examples compare interfaces, inheritance, composition and object creation. They are separate from the patterns confirmed in the restaurant application.

\section{Findings, limitations and references}

The baseline supports two Facade instances and one Observer-like event collaboration. The evolution study introduces online ordering, kitchen station events and a Strategy-based payment structure.

Remaining limitations include incomplete full-stack integration testing, simulated payment success, missing cyclomatic-complexity measurements, incomplete behavior-preservation evidence and incomplete LLM experiment records.

Any LLM-assisted experiment should include the actual model and version, access tier, date, exact prompt, relevant output, source commit and human verification. Missing details must not be invented.

\textbf{Before submission:} verify the baseline commit and license; record source line ranges for the pattern classifications; reconcile the uncertain Chain of Responsibility candidate across all records; commit the evolved application; retain accurate test and complexity limitations; and verify team contributions.

\begin{enumerate}[leftmargin=*,itemsep=1pt]
\item Gamma, E., Helm, R., Johnson, R. and Vlissides, J.
\textit{Design Patterns: Elements of Reusable Object-Oriented Software}.
Addison-Wesley, 1994.
\item Original application repository:
\url{https://github.com/iono-such-things/restaurant-pos-system}.
\item Team repository:
\url{https://github.com/SrivarmaBattini/DP_Project}.
\item React documentation: \url{https://react.dev/}.
\item Express documentation: \url{https://expressjs.com/}.
\item Prisma documentation: \url{https://www.prisma.io/docs}.
\item Socket.IO documentation: \url{https://socket.io/docs/v4/}.
\end{enumerate}

\end{document}
